\section{Conclusions}

Alternating bilateral auditory stimulation should not be bundled into the main ARC-ADHD efficacy claim. It is better treated as a separate proof-of-mechanism question with an acoustically matched non-alternating control.

The key falsification criterion is simple: if left-right alternation does not outperform matched ordinary auditory stimulation on a preregistered proximal recovery outcome, the data do not support a specific bilateral mechanism. This design protects the broader ARC-ADHD programme from overinterpreting an experimental sensory component.
